\chapter{推理时计算缩放——用更多思考换更好答案}
\label{ch:ttc}
%% ============================================================

传统的 LLM 使用方式是：
输入一个问题，模型立刻给出答案。
但 OpenAI 的 o1 系列和 DeepSeek-R1 开创了一种新范式：
\textbf{让模型在回答之前先「思考」}——
生成大量的推理过程（thinking tokens），
然后才给出最终答案。

\section{核心洞察}

\begin{corebox}[推理时计算的颠覆性影响]
作为一个粗略的估计：
\textbf{7B 模型 + 100$\times$ 推理计算 $\approx$ 70B 模型标准推理}
（基于 Snell et al.\ (2024)~\cite{snell2024scaling} 的研究）。
这意味着你可以用一个便宜的小模型，
通过投入更多推理计算来获得与大模型相当的质量。
这彻底改变了「模型大小 = 能力」的等式。
\end{corebox}

传统的范式是\textbf{训练时缩放}（train-time scaling）：
投入更多计算来训练更大的模型，
推理时每个问题使用固定的计算量。
Scaling Laws 描述的就是这种范式。

推理时计算缩放代表了一种\textbf{互补的范式}：
模型大小固定，
但根据问题的难度动态调整推理时的计算量。
简单问题快速回答（几个 token 的思考），
复杂问题深度推理（数千个 token 的思考）。

这两种范式的关系可以类比为人类学习和思考的区别：
训练时缩放相当于\textbf{学更多年}（更多教育），
推理时缩放相当于\textbf{想更久}（更深思考）。
一个博士可能比一个本科生知道更多（训练时缩放），
但本科生如果花足够多的时间仔细思考，
在某些具体问题上也可能得到更好的答案（推理时缩放）。

\begin{whybox}[为什么推理时缩放可能比训练时缩放更高效？]
Snell et al.\ (2024) 的研究给出了一个关键发现：
对于许多具体问题，
将计算预算花在推理阶段比花在扩大模型更具\textbf{计算效率}。

直觉解释：
训练时扩大模型，所有问题都受益于更多的参数——
但简单问题本来就能做对，额外的参数是浪费。
推理时缩放允许\textbf{按需分配}——
简单问题快速回答节省计算，
难的问题多花时间保证质量。
这种动态分配本质上是一种更优的计算资源配置。

类比：
训练时缩放像「给所有学生都上博士课程」——
对已经懂的学生是浪费；
推理时缩放像「允许每个学生根据题目难度决定思考时间」——
更灵活，更高效。
\end{whybox}


\section{推理时缩放的数学框架}

\subsection{Snell et al.\ 的计算最优缩放理论}

Snell et al.~\cite{snell2024scaling} 的核心贡献是建立了
推理时计算的\textbf{缩放法则}（scaling law），
回答了一个关键问题：
给定固定的推理计算预算 $C$，
如何\textbf{最优分配}来最大化性能？

他们分析了两种缩放机制：

\textbf{机制一：基于验证器的搜索}。
生成 $N$ 个候选回答，
用 Process Reward Model（PRM）评分，
选择最优。
性能提升可以建模为：
\begin{equation}
  P(\text{correct} \mid N, \text{PRM}) = f(N, q_{\text{PRM}})
\end{equation}
其中 $q_{\text{PRM}}$ 是验证器的质量，
$N$ 是候选数量。
关键发现是：\textbf{PRM 质量比候选数量更重要}。
一个好的 PRM + 少量候选（$N = 4$）
通常优于一个差的 PRM + 大量候选（$N = 256$）。

\textbf{机制二：自适应分布更新}。
根据具体问题动态调整模型的回答策略。
对于推理模型，这体现为
根据问题难度生成不同长度的推理链。

\textbf{最优分配}的核心结论是：
\begin{equation}
  \text{Optimal strategy}(C, d) = \begin{cases}
    \text{少量深度推理} & \text{如果 } d \text{ 高（难题）} \\
    \text{多次浅层尝试} & \text{如果 } d \text{ 低（简单题）} \\
    \text{放弃推理时缩放} & \text{如果 } d \text{ 极低或极高}
  \end{cases}
\end{equation}
其中 $d$ 是问题的难度。
这个结论有重要的实践意义：
\begin{itemize}[noitemsep]
  \item 对于「模型本来就能做对」的简单问题，
        额外的推理计算几乎没有收益——直接回答即可。
  \item 对于「中等难度」的问题，推理时缩放的收益最大——
        相当于用计算换取了更大模型的效果。
  \item 对于「远超模型能力」的极难问题，
        再多的推理计算也无济于事——
        模型不具备必要的知识或推理能力。
\end{itemize}

\begin{whybox}[推理时缩放的数学等价性]
Snell et al.\ 的一个惊人发现是\textbf{计算等价性}：
一个较小的模型（如 Llama 3.1 8B）
通过最优的推理时计算分配，
可以在某些问题上匹配甚至超过
一个大 14$\times$ 的模型（Llama 3.1 70B）的性能。

形式化地说，对于难度为 $d$ 的问题集合，
存在一个推理计算预算 $C^*(d)$ 使得：
\begin{equation}
  \text{Acc}(\text{小模型}, C^*(d)) \geq \text{Acc}(\text{大模型}, C_0)
\end{equation}
其中 $C_0$ 是标准推理的基线计算量。
当 $C^*(d) < C_{\text{train}}$（训练大模型的额外成本摊到推理上）时，
推理时缩放就是更经济的选择。

但这个等价性不是\textbf{普适}的——
它高度依赖问题难度。
对于极难的问题，没有任何推理计算量
能让小模型匹配大模型。
这就是为什么推理时缩放是训练时缩放的\textbf{互补}而非\textbf{替代}。
\end{whybox}

\subsection{三种 Scaling Laws 的统一视角}

到 2026 年，LLM 领域已经形成了
\textbf{三种 scaling laws} 的统一框架：

\begin{enumerate}[noitemsep]
  \item \textbf{预训练缩放}（Kaplan 2020, Chinchilla 2022）：
        更多的参数 $\times$ 更多的数据 $\times$ 更多的训练计算
        → 更好的基础能力。
        描述的是\textbf{模型知识}的缩放。
  \item \textbf{后训练缩放}（RL、RLHF、SFT）：
        更多的对齐训练 → 更好地遵循指令和推理。
        描述的是\textbf{模型行为}的缩放。
  \item \textbf{推理时缩放}（Snell 2024, o1/R1）：
        更多的推理计算 → 在具体问题上更好的表现。
        描述的是\textbf{问题求解}的缩放。
\end{enumerate}

这三种缩放在\textbf{效率前沿}上各有优势：
预训练缩放适合需要广泛知识的场景，
后训练缩放适合需要特定行为模式的场景，
推理时缩放适合需要深度推理的特定问题。
2026 年的趋势是同时优化这三个维度——
所谓的\textbf{效率缩放}（efficiency scaling），
目标是用 \$1M 的总成本达到以前需要 \$100M 的效果。

\textbf{Budget forcing}（预算强制）
是推理时缩放的一个实用技术：
根据问题的难度\textbf{预先分配}不同的计算预算。
简单问题限制 thinking tokens 在 256 以内，
中等问题分配 2K--8K，
困难问题分配 16K--64K。
这避免了模型在简单问题上「过度思考」浪费资源，
也保证了困难问题有足够的推理深度。


\section{三种策略}

推理时计算缩放有三种基本策略，
各自适用于不同的场景。

\subsection{策略一：扩展推理链（Sequential Scaling）}

\textbf{顺序缩放}让模型在内部进行长链推理——
分解问题、逐步推导、自我验证、纠正错误。
DeepSeek-R1~\cite{deepseekr1} 的 thinking tokens 就是这种策略。
每个 thinking token 都是在「推理计算」上的投资。
模型可能生成数千个 thinking token，
但最终答案的质量远高于直接回答。

一个典型的推理过程可能是这样的：
模型先分析问题结构，
然后尝试一种解法，
发现行不通后自我纠正（``wait, that's not right...''），
换一种方法，
验证中间结果，
最终得到答案。
这种\textbf{自我对话和纠错能力}
是 RL 训练的产物，不是人类预先编程的。

\textbf{思维链推理}（Chain-of-Thought, CoT）
是顺序缩放的基础形式。
即使不是专门训练的推理模型，
简单地在 prompt 中加上
``Let's think step by step'' 就能提升推理质量。
但专门训练的推理模型（如 R1、o1）
的推理深度和自我纠错能力远超简单的 CoT prompting。

\textbf{训练推理模型}的核心方法是
GRPO~\cite{shao2024grpo}（Group Relative Policy Optimization）
配合\textbf{结果导向奖励}（outcome-based reward）。
训练过程不告诉模型「如何推理」，
只告诉它「答案对不对」——
模型必须自己发现有效的推理策略来最大化奖励。
这种纯粹由 RL 驱动的训练方式解释了为什么
R1 会自发地「发明」推理策略：
它不是在模仿人类的推理范例，
而是在\textbf{发现}最大化奖励的行为模式。

顺序缩放的一个微妙之处是\textbf{推理长度与质量}的关系。
并非越长的推理越好——
模型有时会陷入无意义的循环重复，
或者在简单问题上过度展开推理链。
DeepSeek-R1 的训练中观察到，
模型在 RL 早期会出现 ``overthinking'' 现象——
推理链越来越长但答案并没有变好。
通过调整奖励函数（对过长的推理施加轻微惩罚），
可以让模型学会在合适的深度停止思考。

\subsection{策略二：并行采样与 Best-of-N（Parallel Scaling）}

\textbf{并行缩放}对同一个问题生成 $N$ 个独立的回答，
然后用多数投票或奖励模型选最好的。
这相当于「三个臭皮匠赛过诸葛亮」——
即使每个回答独立看不够好，
通过聚合多个独立的回答可以得到更可靠的结果。

\subsubsection{Best-of-N 采样}

\textbf{Best-of-N 采样}是最简单的并行缩放实现：
生成 $N$ 个回答，
用奖励模型（Reward Model, RM）选择分数最高的那个。

数学上，Best-of-N 的效果可以用\textbf{序统计量}分析。
设每个回答的奖励分数 $r_i$ 独立同分布于某分布 $F(r)$。
$N$ 个样本的最大值 $r_{(N)}$ 的期望为：
\begin{equation}
  \mathbb{E}[r_{(N)}] = \int_{-\infty}^{\infty} r \cdot N \cdot F(r)^{N-1} \cdot f(r)\, dr
\end{equation}
当 $F$ 是正态分布 $\mathcal{N}(\mu, \sigma^2)$ 时，
$\mathbb{E}[r_{(N)}]$ 近似为：
\begin{equation}
  \mathbb{E}[r_{(N)}] \approx \mu + \sigma \cdot \Phi^{-1}\!\left(\frac{N}{N+1}\right)
\end{equation}
其中 $\Phi^{-1}$ 是标准正态分布的分位数函数。

这意味着 Best-of-N 的收益是\textbf{对数缩放}的：
从 $N=1$ 到 $N=4$，质量提升显著；
从 $N=16$ 到 $N=64$，提升已经很小；
到 $N=256$，几乎无法进一步提升。
计算成本线性增长但收益对数递减，
所以实践中 $N = 4$--$16$ 是最优的性价比区间。

\begin{warnbox}[Best-of-N 的陷阱：奖励黑客]
Best-of-N 的质量\textbf{严重依赖}奖励模型的质量。
如果奖励模型是训练目标的不完美代理
（这几乎总是如此），
过大的 $N$ 会导致\textbf{奖励黑客}（reward hacking）——
模型生成的回答获得了高奖励分数，
但实际质量并不好，甚至更差。
这是因为大 $N$ 给了模型更多机会
发现奖励模型的弱点并加以利用。

\textbf{正则化 Best-of-N}（RBoN）~\cite{ichihara2025rbon}
通过在目标函数中加入正则项来缓解这个问题：
\begin{equation}
  x^* = \arg\max_{x \in \{x_1, \ldots, x_N\}} \left[ R(x) - \beta \cdot \mathrm{KL}(q \| p_{\text{ref}}) \right]
\end{equation}
正则项惩罚偏离参考分布过远的回答，
防止奖励黑客。
\end{warnbox}

\subsubsection{多数投票与自一致性}

\textbf{多数投票}（Majority Voting / Self-Consistency）~\cite{wang2023selfconsistency}
更适合有\textbf{可验证答案}的问题（如数学题）：
生成 $N$ 个推理过程和答案，
选择出现次数最多的答案。
直觉是：即使每次推理的中间步骤不同，
正确答案更有可能被多次「发现」。

数学分析很清晰：
如果单次回答正确的概率是 $p$，
$N$ 次独立采样中至少有一次正确的概率是：
\begin{equation}
  P(\text{at least one correct}) = 1 - (1-p)^N
\end{equation}
当 $p = 0.3$（单次正确率 30\%）时，
$N = 10$ 次采样就有 $1 - 0.7^{10} \approx 97\%$ 的概率
至少得到一个正确答案。
但从 $N = 10$ 到 $N = 100$，
概率只从 97\% 提升到 99.99\%——
\textbf{收益急剧递减}。

多数投票比 Best-of-N 更\textbf{鲁棒}，
因为它不依赖可能有偏的奖励模型——
它依赖的是「正确答案比错误答案更一致」这个假设。
但它只适用于答案可以\textbf{精确比较}的场景
（数学、选择题、代码输出），
不适用于开放式文本生成。

\subsubsection{奖励模型引导搜索}

更高级的并行缩放策略结合了
奖励模型的评估能力和搜索算法。

\textbf{Weighted majority voting}：
不是简单地数每个答案出现的次数，
而是用奖励模型对每个推理链打分，
然后按分数加权投票。
这结合了自一致性（答案频率）
和质量评估（奖励分数）的优点。

\textbf{LE-MCTS}（Language model Ensemble with MCTS）~\cite{park2025lemcts}：
将多个 LLM 的推理步骤视为
MCTS 中不同的行动选择。
在每一步推理中，选择哪个 LLM 来生成下一步
由 UCB（Upper Confidence Bound）策略决定，
由 PRM 提供每一步的评分。
这实现了\textbf{步骤级别}的模型集成，
比简单的输出级别集成更加精细。

并行缩放的优点是实现极其简单、可线性扩展、
而且每个独立样本可以并行生成，延迟不增加（只增加吞吐需求）。
缺点是效率低——大部分生成是浪费的，
而且对于需要深度推理的问题，
100 次浅层思考可能不如 1 次深度思考。

\subsection{策略三：搜索（Search / Tree of Thought）}

\textbf{搜索}将推理过程视为一棵\textbf{树}，
每个节点是一个中间推理步骤，
每条边是一步推理。
通过系统性地探索多条推理路径，
可以避免单条路径的死胡同。

\textbf{Tree of Thought}（ToT）
是最直观的实现：
在推理的每个关键步骤，
生成多个候选的下一步，
评估每个候选的质量，
选择最有希望的继续展开。

\textbf{MCTS}（Monte Carlo Tree Search）
将 AlphaGo 的搜索算法引入语言推理：
从根节点（问题）出发，
通过选择（Selection）、扩展（Expansion）、
模拟（Simulation）、回传（Backpropagation）
四个步骤反复迭代，
逐步构建推理树并找到最优路径。

搜索的关键组件是\textbf{评估函数}——
如何判断一个中间推理步骤的质量？
这里引入了 \textbf{Process Reward Model}（PRM）的概念：
与传统的 Outcome Reward Model（ORM，只评分最终答案）不同，
PRM 对\textbf{每个推理步骤}打分。
这提供了更细粒度的信号——
不仅知道最终答案对不对，
还知道从哪一步开始出了问题。
PRM 使得搜索可以\textbf{及早剪枝}错误的推理路径，
大幅提高搜索效率。

\begin{whybox}[ORM vs PRM：粗粒度 vs 细粒度奖励]
考虑一个数学问题的 5 步推理过程，
其中第 3 步出了错误。

\textbf{ORM}只看最终答案：答案错了 → 整个推理链得低分。
但 ORM 不知道前 2 步是对的，
导致在搜索中可能丢弃那些「前半部分优秀」的推理路径。

\textbf{PRM}对每一步评分：
Step 1: 0.95, Step 2: 0.92, Step 3: 0.15, Step 4: 0.10, Step 5: 0.08。
PRM 清晰地定位了错误在 Step 3，
搜索算法可以保留 Step 1--2 的推理，
只从 Step 3 开始尝试新的方向。

PRM 的训练更困难——
需要\textbf{步骤级别}的标注，
而不是简单的「答案对不对」。
但其提供的细粒度信号
使得搜索效率提升数个数量级。
Snell et al.\ 的实验表明，
PRM 引导的 beam search 在相同计算预算下
比 ORM 引导的 Best-of-N 高出 10--20\% 的准确率。
\end{whybox}

\textbf{AlphaProof} 和 \textbf{AlphaGeometry 2}
是搜索策略的极致展示：
在 2024 年国际数学奥林匹克（IMO）中，
AlphaProof 与 AlphaGeometry 2 联合解出了 6 道题中的 4 道
（其中 AlphaProof 解决了代数和数论题，
AlphaGeometry 2 解决了几何题），
达到了银牌水平。
AlphaProof 使用 Gemini 模型将自然语言转换为 Lean 4 形式化语言，
再通过类似 AlphaZero 的强化学习方法进行证明搜索。
每一步都可以被\textbf{机器验证}（proof checker），
消除了幻觉的可能性。

\textbf{Beam search over reasoning chains}
是一种更轻量的搜索方法：
类似于传统的 beam search，
维护 $B$ 条最有希望的推理链，
每步扩展每条链，保留总分最高的 $B$ 条。
PRM 提供每一步的评分。
这比完整的 MCTS 计算量更小，
但仍然比纯顺序推理有显著的质量提升。

\subsection{三种策略的对比}

\begin{keybox}[推理时计算策略对比]
\begin{itemize}[noitemsep]
  \item \textbf{顺序缩放}：一次深度思考。
        优势——推理深度高，适合需要多步推导的问题。
        劣势——单次推理可能走进死胡同。
  \item \textbf{并行缩放}：$N$ 次独立思考。
        优势——实现简单，可并行，延迟不增加。
        劣势——效率低，收益递减，可能遭受奖励黑客。
  \item \textbf{搜索}：系统性探索推理树。
        优势——可以回溯，避免死胡同，理论上最优。
        劣势——需要 PRM，计算量大，实现复杂。
\end{itemize}
经验法则：低难度问题用并行缩放，
中高难度用顺序缩放，
极高难度（数学竞赛、形式化证明）用搜索。
\end{keybox}

\subsection{混合缩放（Hybrid Scaling）}

\textbf{混合缩放}结合并行和顺序——
并行生成多个推理链，
每条链内部是顺序推理，最后取最优。
ThreadWeaver（2025 年提出）就是这种方法，
实现了 $1.53\times$ 的延迟降低同时保持准确率。

实践中最有效的配置通常是混合的：
让推理模型生成 $N$（通常 4--16）个独立的长推理链，
然后用多数投票或奖励模型选最优。
这同时利用了「深度思考」和「多样性」的优势。


\section{DeepSeek-R1：从零开始学会推理}

DeepSeek-R1~\cite{deepseekr1} 是推理时缩放领域的里程碑。
它首次证明了\textbf{纯粹的强化学习}（无需人类推理示范）
就能让模型自发地发展出复杂的推理能力。

2025 年 9 月，R1 论文以封面文章发表于
\textit{Nature}（Vol.~645, Issue~8081, 2025 年 9 月 17 日），
成为\textbf{首篇通过 Nature 独立同行评审的 LLM 论文}。
梁文锋（Liang Wenfeng）为通讯作者。
Nature 编辑部以``AI can learn to show its workings through trial and error''为题
发表评论，确认了 R1 的核心发现：
纯 RL（无需人类标注的推理轨迹）能够激励高级推理行为的涌现——
包括自我反思（self-reflection）、
验证（verification）和动态策略调整（dynamic strategy adaptation）。
这一里程碑标志着 LLM 推理研究
从工程实践迈入了经由严格同行评审认可的科学领域。

\subsection{R1-Zero：从 RL 中涌现的推理}

在 R1 的完整训练之前，
DeepSeek 团队先进行了一个大胆的实验：
在 DeepSeek-V3-Base 上\textbf{直接}应用 GRPO RL 训练，
\textbf{不使用任何 SFT 数据}。
这个模型被称为 \textbf{DeepSeek-R1-Zero}。

R1-Zero 的训练设计极其简洁：
\begin{itemize}[noitemsep]
  \item \textbf{奖励信号}：只有两类——
        数学题答案是否正确（accuracy reward），
        和输出格式是否符合要求（format reward）。
        没有过程奖励（PRM），没有人类偏好。
  \item \textbf{训练数据}：数学、代码和逻辑推理题，
        全部有\textbf{可验证}的正确答案。
  \item \textbf{优化算法}：GRPO——
        对每个问题生成一组回答，
        根据组内相对排名计算优势函数，
        不需要单独的价值网络或奖励模型。
\end{itemize}

\subsection{``Aha Moment''——模型的顿悟时刻}

R1-Zero 训练过程中出现了令研究者震惊的现象。
在 RL 训练到一定阶段后，
模型自发地开始展现出一系列
\textbf{从未在训练数据中出现过}的推理行为：

\begin{corebox}[DeepSeek-R1 的 ``Aha Moment'']
在训练中期的某次评估中，
模型在解数学题时写下了这样的推理过程：

\textit{``Wait, wait. Wait. That's an aha moment.
I can actually solve this by re-reading the problem
and approaching it differently...
Let me reconsider what I assumed earlier...''}

这不是人类编写的推理模板——
模型自己「发明」了自我反思和重新审视的策略。
这个时刻被 DeepSeek 团队称为 ``aha moment''，
因为它标志着模型从简单的模式匹配
跃迁到了\textbf{元认知}层面的推理。
\end{corebox}

具体而言，R1-Zero 自发发展出了以下推理策略：

\textbf{回溯}（Backtracking）：
模型在推理链中发现前面的方向不对时，
主动放弃当前路径，回退到更早的步骤重新推导。
这是一种典型的搜索行为——
模型学会了在推理树中「剪枝」和「回溯」。

\textbf{自我验证}（Self-verification）：
模型在得出中间结论后，
会主动检验该结论的正确性——
例如将答案代回原方程验证，
或用不同方法重新推导检查一致性。
这模拟了数学家的证明检查过程。

\textbf{分类讨论}（Case analysis）：
面对多种可能情况的问题，
模型学会了系统性地枚举每种情况，
逐一分析，最后综合得出结论。

\textbf{反证法}（Proof by contradiction）：
在某些逻辑推理问题中，
模型会假设结论不成立，
推导出矛盾，从而证明结论成立。

这些策略的涌现具有深远的意义：
它表明复杂的推理能力\textbf{不需要}显式的人类教学，
可以通过适当的 RL 训练从模型中自然涌现。
这与 AlphaGo Zero 的启示一致——
在有可验证奖励信号的领域，
RL 能够发现人类未曾想到的策略。

\subsection{R1 的完整训练流程}

R1-Zero 虽然展现了令人兴奋的推理涌现，
但存在两个严重的问题：
\begin{itemize}[noitemsep]
  \item \textbf{可读性差}：输出混合多种语言，
        推理链混乱难读，对用户不友好。
  \item \textbf{非推理任务退化}：
        在日常对话、翻译等不需要推理的任务上
        表现明显下降。
\end{itemize}

DeepSeek-R1 的完整训练分为四个阶段：

\textbf{阶段一：冷启动 SFT}。
收集少量高质量的推理示范数据（数千条），
对基础模型做 SFT。
这一步不是为了教模型推理——
而是为了稳定后续 RL 训练的起点，
避免 R1-Zero 中的格式混乱问题。

\textbf{阶段二：推理导向 RL}。
在冷启动模型上做 GRPO RL 训练，
使用数学、代码和逻辑题的可验证奖励。
这是推理能力形成的核心阶段。

\textbf{阶段三：拒绝采样 + 通用 SFT}。
用阶段二的模型生成大量推理链，
通过拒绝采样保留高质量的输出（约 800K 条），
混合通用的 SFT 数据（日常对话、翻译等），
做一轮全面的 SFT。
这一步恢复了模型在非推理任务上的能力。

\textbf{阶段四：全场景 RL}。
在所有类型的任务上做最终的 RL 训练，
使用混合的奖励信号
（可验证任务用规则奖励，开放式任务用奖励模型），
进一步提升整体质量和对齐度。

\subsection{R1 vs o1/o3：两条不同的路径}

DeepSeek-R1 和 OpenAI 的 o1/o3 系列
代表了推理时缩放的两条不同路径：

\begin{table}[H]
\centering\small
\caption{DeepSeek-R1 vs OpenAI o 系列对比}
\begin{tabularx}{\linewidth}{l X X}
\toprule
\textbf{维度} & \textbf{DeepSeek-R1} & \textbf{OpenAI o1/o3} \\
\midrule
\textbf{核心方法} & GRPO RL + 可验证奖励。不需要单独的奖励模型。 & RL + 详细的 PRM。使用复杂的奖励建模系统。 \\
\textbf{推理过程} & 完全可见。用户可以看到完整的 thinking tokens。 & 隐藏。用户只看到最终答案（部分摘要可见）。 \\
\textbf{开源} & 模型权重完全开源，训练方法论在论文中详细公开。 & 完全闭源，仅通过 API 提供服务。 \\
\textbf{基础架构} & 671B MoE（37B 活跃参数），基于 DeepSeek-V3。 & 未公开。推测基于 GPT-4 级别架构。 \\
\textbf{蒸馏} & 提供 1.5B 到 70B 的完整蒸馏系列。 & 无官方蒸馏，o3-mini 作为轻量版本。 \\
\textbf{性能（AIME 2024）} & 79.8\% & o1: 83.3\%, o3-mini (high): 87.3\% \\
\bottomrule
\end{tabularx}
\end{table}

o3-mini 是一个值得关注的案例：
它在多个数学基准（如 AIME）上\textbf{超过}了 o1，
同时成本只有 o1 的 $1/14$（\$1.10/\$4.40 vs \$15/\$60 per M tokens）。
这证明了推理时缩放中
\textbf{算法效率的提升}可以同时提高质量和降低成本。

2025 年 4 月，OpenAI 发布了 o3 和 o4-mini，
标志着推理模型的又一次飞跃。
o3 和 o4-mini 引入了\textbf{工具使用}能力——
它们可以在推理过程中自主决定
搜索网页、执行代码、分析图像，
将工具调用和推理无缝结合。
这代表了从「纯思考」到「思考 + 行动」的范式转变。

\subsection{推理基准的军备竞赛（2025--2026）}

DeepSeek-R1 发表后，推理模型的竞争进入白热化阶段。
2025 年末至 2026 年初，多个模型在传统推理基准上
取得了接近天花板的成绩：

\begin{table}[H]
\centering\small
\caption{推理模型基准成绩（2025 年末至 2026 年初）}
\begin{tabularx}{\linewidth}{l X r}
\toprule
\textbf{模型} & \textbf{关键成绩} & \textbf{发布时间} \\
\midrule
\textbf{DeepSeek V3.2-Speciale} &
  IMO 金牌（35/42）、ICPC 第 2 名（492/600）、IOI 第 10 名、AIME 96.0\%（超越 GPT-5-High 94.6\%） &
  2025-12 \\
\textbf{Doubao Seed 2.0 Pro} &
  AIME 2025 98.3\%、Codeforces Rating 3020 &
  2026-02 \\
\textbf{Step 3.5 Flash} &
  AIME 97.3\%，仅 11B 活跃参数（196B 总参数 MoE） &
  2026-02 \\
\textbf{Kimi K2.5} &
  AIME 96.1\%、HMMT 95.4\%、GPQA Diamond 87.6\% &
  2026-01 \\
\textbf{Qwen3-Max-Thinking} &
  AIME/HMMT 部分基准宣称 100\% &
  2025-09+ \\
\bottomrule
\end{tabularx}
\end{table}

值得注意的是 \textbf{DeepSeek V3.2-Speciale} 的竞赛表现：
IMO 金牌（35/42 分）意味着模型的数学推理能力
已经达到了国际数学奥林匹克顶尖选手的水平，
而 ICPC 第 2 名和 IOI 第 10 名
则证明了这种推理能力可以迁移到编程竞赛领域。
\textbf{Step 3.5 Flash} 尤其引人关注——
它仅用 11B 活跃参数就达到了 AIME 97.3\%，
证明了小型 MoE 架构配合高质量推理训练
可以在效率和性能之间取得惊人的平衡。

\begin{warnbox}[基准饱和与评估范式转移]
多个模型在 AIME 上突破 96\% 甚至宣称 100\%，
这种\textbf{分数向天花板收敛}的现象
表明传统推理基准正在快速饱和。
当区分度消失时，基准就失去了评估价值。

下一代有意义的评估正在转向更接近真实世界的任务：
\textbf{SWE-bench Pro}（1,865 个多语言真实代码修复任务，
当前最高约 44\%）衡量的是工程推理能力；
\textbf{BrowseComp}（Web agent 任务）
衡量的是模型在开放环境中的规划和执行能力。
这些基准的天花板远高于 AIME，
更能区分不同模型的实际推理水平。
\end{warnbox}


\section{工具增强推理}

2025--2026 年，推理模型的一个重要演进方向是
将\textbf{外部工具}集成到推理过程中。
模型不再仅仅依赖内部推理——
而是可以在思考过程中调用工具来辅助推理。

\subsection{代码执行作为推理工具}

最自然的工具增强推理是\textbf{代码执行}。
在解决数学或数据分析问题时，
模型可以编写代码来精确计算中间结果，
而不是依赖自然语言推理（容易出现算术错误）。

典型的流程是：
\begin{enumerate}[noitemsep]
  \item 模型分析问题，决定使用代码来辅助计算。
  \item 在 thinking 过程中生成 Python 代码。
  \item 代码在沙箱环境中执行，返回结果。
  \item 模型将执行结果纳入后续推理链。
\end{enumerate}

这解决了 LLM 推理中最臭名昭著的弱点之一：
\textbf{算术错误}。
LLM 通过自然语言做多步算术时，
每一步都有累积误差的风险。
而代码执行保证了计算的\textbf{精确性}。
OpenAI 的 o3/o4-mini 原生支持
在推理过程中执行 Python 代码——
这使得它们在需要精确计算的数学和科学问题上
表现显著优于仅用自然语言推理的模型。

\subsection{搜索与检索增强推理}

推理过程中的\textbf{检索增强}（Retrieval-Augmented Reasoning）
让模型在思考时可以搜索外部知识。
当模型在推理中遇到不确定的事实性问题时，
它可以暂停推理、搜索相关信息、
然后基于检索到的事实继续推理。

这与传统的 RAG（Retrieval-Augmented Generation）不同：
\begin{itemize}[noitemsep]
  \item \textbf{传统 RAG}：在生成\textbf{之前}检索一次，
        将检索结果作为上下文输入模型。
  \item \textbf{推理中检索}：在推理\textbf{过程中}多次检索，
        每次检索都基于当前推理状态的需要。
        检索和推理交替进行，形成「思考 → 搜索 → 思考」的循环。
\end{itemize}

这种交互式检索在需要综合多源信息的推理任务中
效果尤为显著——
例如法律案例分析（需要检索多个先例）
或科学文献综述（需要交叉引用多篇论文）。

\subsection{多步 Agent 推理}

\textbf{Agent 推理}是工具增强推理的最完整形式。
推理模型不仅可以思考和使用单个工具，
还可以\textbf{规划}和\textbf{编排}多步骤的工作流。

Claude Opus 4.6 的 extended thinking
允许模型在一个统一的推理过程中
思考、调用工具、分析工具返回的结果、
再继续思考——
整个过程是一个\textbf{连贯的}推理链，
而不是分离的「思考步骤」和「工具调用步骤」。

OctoTools~\cite{lu2025octotools} 提出了一个
可扩展的 agent 推理框架。
通过\textbf{工具卡片}（tool card）机制，
模型可以在无需额外训练的情况下
集成新的工具——
每张工具卡片描述了工具的功能、输入/输出格式
和使用时机。
模型通过阅读工具卡片来决定何时使用什么工具。
实验表明，即使是最强的 LLM + 完善的 agent 工作流，
在复杂的跨论文科学推理任务上
准确率也仅有 38.78\%，
困难子集上更是低至 18.47\%——
这表明工具增强推理仍有巨大的改进空间。

\begin{keybox}[工具增强推理的层次]
\begin{itemize}[noitemsep]
  \item \textbf{Level 0——纯推理}：
        模型仅用自然语言推理，不使用任何外部工具。
        CoT、R1-Zero 属于这一层次。
  \item \textbf{Level 1——单工具推理}：
        模型在推理中使用单一工具（如代码执行器）。
        o3 的 Python 执行属于这一层次。
  \item \textbf{Level 2——多工具推理}：
        模型在推理中编排多个工具（代码、搜索、API 调用）。
        o3/o4-mini 的完整工具使用属于这一层次。
  \item \textbf{Level 3——Agent 推理}：
        模型规划多步骤工作流，
        包括子任务分解、工具编排、结果验证、
        和基于中间结果的动态调整。
        Claude 的 extended thinking + tool use 属于这一层次。
\end{itemize}
\end{keybox}


\section{推理模型的部署经济学}

推理时计算缩放正在深刻重塑 AI 的经济学。
推理需求预计将大幅超过训练需求（据行业估计可达百倍量级）。
GPT-4 等效性能的成本从 2022 年的 \$20/M tokens
降至 2026 年的 \$0.40/M tokens——
$1000\times$ 的下降，主要来自推理优化。

但推理模型也带来了新的成本挑战：
DeepSeek-R1 在回答一个数学问题时
可能生成 10,000+ 个 thinking token，
然后才给出 100 个 token 的答案。
thinking token 虽然不展示给用户，
但同样需要计算和内存资源。

\subsection{何时使用推理时缩放}

推理时缩放不是万灵药——
对于不同类型的任务，性价比差异巨大：

\textbf{高价值场景}（值得投入 10--100$\times$ 推理计算）：
\begin{itemize}[noitemsep]
  \item 数学证明和竞赛题——答案可验证，推理深度直接决定正确率。
  \item 代码 bug 修复和系统设计——错误成本高，
        多花 \$0.10 的推理费用可以避免 \$10,000 的生产事故。
  \item 法律和医疗分析——高风险决策，准确性比延迟更重要。
  \item 科学研究中的假设生成——需要深度推理和创造性思考。
\end{itemize}

\textbf{低价值场景}（标准推理即可）：
\begin{itemize}[noitemsep]
  \item 日常聊天和问答——首次回答通常就足够好。
  \item 翻译和摘要——任务结构明确，不需要深度推理。
  \item 简单的分类和提取——模型直觉足够准确。
\end{itemize}

\subsection{路由与预算管理}

这催生了\textbf{推理预算管理}的概念：
对于简单问题，限制 thinking token 的数量以节省成本；
对于困难问题，分配更多的推理预算以保证质量。

\textbf{难度路由}（Difficulty Routing）
是一种经济高效的方案：
使用一个轻量级的分类器（甚至是 prompt-based 的判断）
评估问题的难度，
将简单问题路由到标准模型（快且便宜），
将困难问题路由到推理模型（慢但准确）。
好的路由器可以将 80\% 的请求用便宜的方式处理，
只对 20\% 真正需要深度推理的问题投入额外计算。

\begin{corebox}[可配置的推理预算：行业实践]
Claude Opus 4.6 支持可配置的 thinking budget，
允许开发者通过 API 设置任意的 \texttt{budget\_tokens} 参数
来控制推理深度。
这将推理时缩放从「全有或全无」
变成了\textbf{连续可调}的维度——
开发者可以根据任务需求和成本预算
在延迟、质量和成本之间做精细的权衡。

Qwen3~\cite{qwen2025qwen3} 走了另一条路：
在单个模型中统一 thinking 和 non-thinking 模式。
用户可以根据需要在两种模式间切换——
简单问题直接回答（non-thinking），
复杂问题开启深度推理（thinking）——
不需要部署两个不同的模型。

Google Gemini 2.5/3 引入了\textbf{动态思考模式}——
模型根据问题复杂度自动调整推理深度，
无需用户手动控制。
Flash 版本提供快速推理，Pro 版本提供深度推理，
两者共享同一架构但分配不同的计算预算。
\end{corebox}

\subsection{共享推理的可能性}

对于多个用户问相似的难题（如热门的面试编程题），
可以缓存和复用推理过程。
一个用户的深度推理结果可以被后续用户共享，
避免重复计算。
这类似于 prefix caching 的思想，
但扩展到了推理链（reasoning chain）层面。
这在教育和客服场景中有巨大的潜力。

\subsection{成本模型}

一个推理模型回答一个数学问题的成本构成：
\begin{itemize}[noitemsep]
  \item \textbf{标准回答}（无 thinking）：
        $\sim$500 输出 token，成本 $\sim$\$0.001。
  \item \textbf{浅层推理}（1K thinking tokens）：
        $\sim$1500 总 token，成本 $\sim$\$0.003。
  \item \textbf{深度推理}（10K thinking tokens）：
        $\sim$10500 总 token，成本 $\sim$\$0.02。
  \item \textbf{Best-of-16 + 深度推理}：
        $\sim$168K 总 token，成本 $\sim$\$0.30。
\end{itemize}

深度推理的成本是标准回答的 20 倍，
Best-of-16 + 深度推理是 300 倍。
但如果问题本身的价值足够高
（一道竞赛题的解答、一个关键 bug 的修复），
这个成本仍然极为合算。

关键在于\textbf{匹配推理投入和问题价值}——
用 \$0.30 的推理计算来回答「今天天气怎么样」是浪费，
用 \$0.001 的快速回答来解数学竞赛题是不足。
好的推理时计算系统需要\textbf{自适应}地分配资源。


\section{推理模型的前沿与未来}

\subsection{从推理到行动：Agent 范式}

2026 年的一个核心趋势是
推理模型从「思考者」进化为「行动者」。
传统推理模型只在\textbf{文本空间}中推理；
新一代模型在\textbf{行动空间}中推理——
它们可以规划行动序列、执行工具调用、
观察结果、调整计划。

这种「推理 + 行动」的范式
本质上是将推理时缩放应用于更广阔的问题空间。
不再只是「想更久」，
而是「想更久，并且在想的过程中做实验来验证想法」。

\subsection{效率缩放：做更多、花更少}

2026 年行业的核心议题不再是「更大的模型」，
而是\textbf{效率缩放}——
用更少的计算达到同等甚至更好的效果。
这体现在多个层面：
\begin{itemize}[noitemsep]
  \item \textbf{蒸馏}：
        R1 的 7B 蒸馏模型在 MATH-500 上达到 92.8\%，
        接近 70B 级别——
        1/10 的参数，1/10 的推理成本。
  \item \textbf{混合推理模式}：
        Qwen3 的 thinking/non-thinking 统一架构
        让一个模型适应所有场景，
        避免了部署两套模型的开销。
  \item \textbf{推理预算自适应}：
        根据问题难度动态分配计算，
        80\% 的简单请求用最少的计算完成。
  \item \textbf{推理链缓存}：
        复用相似问题的推理过程，
        避免重复的深度推理计算。
\end{itemize}

\subsection{开放的问题}

推理时缩放仍有多个\textbf{未解决的问题}：

\textbf{如何训练更好的 PRM？}
PRM 的质量决定了搜索策略的上限，
但获取步骤级别的训练数据极为昂贵。
自动化的 PRM 训练（如通过 MCTS 回溯自动标注）
是一个活跃的研究方向。

\textbf{推理时缩放有没有「上限」？}
Snell et al.\ 表明，在极难的问题上，
增加推理计算的收益会趋近于零。
但目前尚不清楚这个上限在哪里，
以及更好的推理策略能否推高这个上限。

\textbf{推理能力能否泛化？}
在数学和代码上训练的推理能力
能否迁移到法律推理、医学诊断或科学发现？
现有证据表明迁移是有限的——
推理的\textbf{形式}可以迁移（如回溯、验证），
但\textbf{领域知识}不能自动获得。

\textbf{什么是推理的本质？}
当前的推理模型是否真的在「推理」，
还是只是在做更复杂的模式匹配？
这个哲学问题对实践也有影响——
如果只是模式匹配，那么在训练分布外的问题上
推理能力会急剧下降。

这些问题的答案将决定
推理时缩放能走多远，
以及它最终能否像预训练缩放那样
带来持续的、可预测的性能提升。


%% ============================================================
%% NEW REFERENCES
%% ============================================================
% NEW REF: wang2023selfconsistency = Wang, Xuezhi, et al. "Self-Consistency Improves Chain of Thought Reasoning in Language Models." ICLR 2023.
% NEW REF: ichihara2025rbon = Ichihara, Yuki, et al. "Evaluation of Best-of-N Sampling Strategies for Language Model Alignment." TMLR 2025.
% NEW REF: park2025lemcts = Park, Sungjin, et al. "Ensembling Large Language Models with Process Reward-Guided Tree Search for Better Complex Reasoning." NAACL 2025.
% NEW REF: lu2025octotools = Lu, Pan, et al. "OctoTools: An Agentic Framework with Extensible Tools for Complex Reasoning." arXiv:2502.11271, 2025.
